%═══════════════════════════════════════════════════════════════════════════════
\section{Conclusion}\label{sec:discussion}
%═══════════════════════════════════════════════════════════════════════════════

\subsection{What the Thesis Shows}\label{sec:disc_answer}

The loss must enter the construction of the evidence, not only its
interpretation. For a named target, the profile proportional to its loss is
the unique safest choice and leads to it. Information can be added only with
the target held as a constraint, since the likelihood ratio can lead
elsewhere even when the belief is right, and no reshaping of the ratio does
this job unless the loss is a function of it. The shopkeeper application
shows that this matters in practice.

When the model and the belief favour different actions, the disagreement
should be reported, since which to follow is the decision-maker's call.
Reporting the likelihood ratio alone, without the loss and the action it
induces, commits silently to an objective the decision may not share.

\subsection{Limitations}\label{sec:disc_limits}

Everything here is proved for finitely many outcomes. Theorem~\ref{thm:cert}
should extend to richer outcome spaces, but the thesis does not prove that,
and the frontier depends more on finiteness: with many actions or a continuum
of outcomes, both the existence of an optimum and the shape of the edges where
the target ties with a competitor need further assumptions.

Proposition~\ref{prop:boundary} also assumes that exactly one validity
constraint binds at the optimum. This holds for the one-point models it is
applied to, but with a composite model two constraints can bind at once, and
the most informative profile is then a compromise between members rather than
a ratio against any single one.

The information criterion is a choice, not a claim about what decision-makers
want. A different measure would move the interior of the frontier but not the
protective endpoint, so the target, and how much protection to trade for
information, must still be named by the decision-maker.

The empirical work is conditional too. The simulations construct their
distributions, so they show how the mechanism behaves rather than rank rules
in general. The application uses a real belief, but its loss table and model
are the analyst's choice. Section~\ref{sec:mkt_shape} varies the input that is
hardest to defend, the historical rate, and the finding survives, but a
different loss table could change it.

\subsection{Open Questions}\label{sec:disc_open}

Three questions remain open. The first is to extend the frontier beyond the
finite setting. The second is to find loss tables that are determined by the
likelihood ratio in practice, and to measure what is lost when they are not.
The third is how the model $\mathcal P$ should be chosen. It sets the strength
of every guarantee here, yet in the application it was assembled from
defensible judgements rather than derived from a principle.
